\chapter{Sleep}

\section*{Definition and Overview}
Sleep concerns lifestyle factors -- diet, physical activity, sleep, substance use, and social engagement -- that influence health and disease. Lifestyle medicine is an evidence-based discipline that addresses these modifiable determinants to prevent, manage, and sometimes reverse chronic disease. The cumulative impact of daily habits on long-term health outcomes is substantial.

\section*{Epidemiology}
Modifiable lifestyle factors are responsible for a large proportion of chronic disease burden globally. Tobacco use remains the leading preventable cause of death. Physical inactivity, unhealthy diet, and harmful alcohol use are major contributors to cardiovascular disease, type 2 diabetes, and several cancers. Obesity prevalence has tripled in many countries since 1975. Mental health and social connection are increasingly recognised as independent determinants of physical health.

\section*{Aetiology and Pathophysiology}
\begin{itemize}
    \item \textbf{Diet:} excess energy intake, high sugar/salt/saturated fat, low fibre, and ultra-processed foods drive obesity, metabolic syndrome, and cardiovascular disease
    \item \textbf{Physical inactivity:} contributes to insulin resistance, sarcopenia, osteoporosis, and cardiovascular deconditioning
    \item \textbf{Sleep disturbance:} chronic sleep deprivation disrupts hormonal regulation (cortisol, ghrelin, leptin) and immune function
    \item \textbf{Substance use:} tobacco, alcohol, and other drugs cause direct tissue damage and increase disease risk
    \item \textbf{Chronic stress:} sustained activation of the HPA axis and sympathetic nervous system
    \item \textbf{Social isolation:} associated with increased all-cause mortality comparable to smoking
\end{itemize}

\section*{Clinical Features}
\begin{itemize}
    \item Fatigue and low energy
    \item Unintended weight gain or loss
    \item Poor sleep quality or insomnia
    \item Reduced exercise tolerance or muscle weakness
    \item Low mood, irritability, or anxiety
    \item Frequent minor illness (immune suppression)
    \item Reduced productivity and social engagement
\end{itemize}

\section*{Differential Diagnosis}
\begin{itemize}
    \item Underlying medical condition causing the symptom (e.g.\ hypothyroidism causing fatigue; anaemia causing tiredness)
    \item Depression or anxiety disorder
    \item Substance use disorder
    \item Eating disorder
    \item Chronic fatigue syndrome / ME
    \item Medication side effects
\end{itemize}

\section*{When to Seek Medical Care}
\begin{itemize}
    \item Before starting a new exercise or diet program if you have a long-term health condition
    \item If you have symptoms such as chest pain, dizziness, or persistent low mood
    \item If you are concerned about your weight, alcohol use, or smoking
    \item If lifestyle changes are not improving symptoms despite sustained effort
\end{itemize}

\section*{Diagnosis and Investigations}
\begin{itemize}
    \item \textbf{Lifestyle assessment:} structured history of diet, activity, sleep, substance use, and social engagement
    \item \textbf{Anthropometry:} weight, height, BMI, waist circumference
    \item \textbf{Blood tests:} lipid profile, HbA1c, liver function, thyroid function, vitamin D
    \item \textbf{Blood pressure:} measurement and monitoring
    \item \textbf{Mental health screening:} depression and anxiety screening questionnaires
    \item \textbf{Sleep assessment:} sleep diary, Epworth Sleepiness Scale, or polysomnography if sleep apnoea suspected
\end{itemize}

\section*{Management}
\begin{itemize}
    \item \textbf{SMART goals:} specific, measurable, achievable, relevant, and time-bound behavioural targets
    \item \textbf{Diet:} Mediterranean or DASH dietary pattern; increase vegetables, fruit, whole grains; reduce ultra-processed foods, added sugar, and salt
    \item \textbf{Physical activity:} build up to 150--300 minutes of moderate-intensity activity per week, plus resistance training twice weekly
    \item \textbf{Sleep:} consistent sleep routine; 7--9 hours per night; address sleep disorders
    \item \textbf{Smoking cessation:} behavioural support plus nicotine replacement therapy, varenicline, or bupropion
    \item \textbf{Alcohol reduction:} brief intervention, motivational interviewing, or specialist referral
    \item \textbf{Stress management:} mindfulness, meditation, cognitive behaviour therapy, or other psychological strategies
\end{itemize}

\section*{Complications}
\begin{itemize}
    \item Cardiovascular disease: coronary artery disease, stroke, hypertension
    \item Metabolic disease: type 2 diabetes, metabolic syndrome, non-alcoholic fatty liver disease
    \item Certain cancers (colorectal, breast, endometrial, kidney)
    \item Mental health disorders: depression, anxiety
    \item Musculoskeletal: osteoarthritis, osteoporosis, sarcopenia
    \item Social and occupational consequences of substance use
\end{itemize}

\section*{Prognosis and Outcomes}
Lifestyle modifications can prevent, delay, or improve many chronic conditions. Even modest weight loss (5--10\%) improves metabolic markers. Regular physical activity reduces all-cause mortality by approximately 30\%. Smoking cessation reduces cardiovascular risk within months. Behavioural change is most sustainable when introduced gradually and supported by social and environmental structures. Relapse is common and should be anticipated and managed without judgement.

\section*{Prevention and Self-Care}
\begin{itemize}
    \item Make one change at a time; set realistic goals
    \item Find physical activities you enjoy so they are easier to sustain
    \item Plan meals and movement into your weekly routine
    \item Prioritise sleep as a foundation of health
    \item Build strong social connections; stay engaged with community
    \item Ask for support from family, friends, a GP, or a structured program
    \item Monitor progress and celebrate improvements
\end{itemize}
